% =====================================================================
%  Phụ lục — Tái lập kết quả
% =====================================================================
\chapter{Tái lập kết quả}
\label{ch:phuluc}

\section{Môi trường}

Windows 11, Python 3.13, GPU NVIDIA GeForce RTX 4060 Laptop cho PyTorch (CUDA 12.6), CPU cho TensorFlow. PyTorch
\SL{bench/torch_version}, Keras \SL{bench/keras_version} với TensorFlow 2.21. Mọi thư viện cài trong môi trường ảo riêng
của bài, liệt kê trong \texttt{requirements.txt}.

\section{Các bước}

\begin{enumerate}
  \item Tạo môi trường và cài thư viện:
\begin{verbatim}
cd "src/Assignment 06"
python -m venv .venv
.venv\Scripts\activate
pip install torch==2.13.0 --index-url https://download.pytorch.org/whl/cu126
pip install -r requirements.txt
python -m ipykernel install --user --name a06
\end{verbatim}
  \item Tải \texttt{train\_v2.csv.7z} và \texttt{user\_logs\_v2.csv.7z} của cuộc thi KKBox trên Kaggle (cần đăng nhập và
        chấp nhận luật cuộc thi) vào \texttt{data/kkbox/}, giải nén. \texttt{AMZN.csv} đã có sẵn trong kho.
  \item Chạy các notebook theo thứ tự số. Notebook 03, 04 (GPU) và 05 (CPU) chạy song song được. Trước khi chạy, đặt
        biến môi trường \texttt{OPENBLAS\_NUM\_THREADS} bằng 8: trên máy 32 luồng, OpenBLAS mặc định cấp bộ đệm cho 32
        luồng và có lúc làm kernel chết vì không cấp phát được bộ nhớ.
\begin{verbatim}
cd notebooks
jupyter nbconvert --to notebook --execute --inplace ^
    --ExecutePreprocessor.kernel_name=a06 --ExecutePreprocessor.timeout=-1 00_*.ipynb
\end{verbatim}
  \item Kiểm chứng Web App: \texttt{python app/smoke\_test.py}; chạy thật: \texttt{python app/app.py}.
  \item Dựng báo cáo: \texttt{python "Report/Assignment 06/build.py"} (cần LuaLaTeX, \texttt{latexmk}, \texttt{biber}
        và gói \texttt{minted}).
\end{enumerate}

Lần chạy đầu của notebook 02 gom log KKBox thành tệp \texttt{kkbox\_seq.npz}; các lần sau đọc thẳng tệp này.

\maNguon{data_churn__build_cache}{Gom dữ liệu KKBox một lần và lưu lại}
